% !TEX program = pdflatex
\documentclass[11pt,a4paper]{article}

\usepackage[utf8]{inputenc}
\usepackage[T1]{fontenc}
\usepackage{amsmath,amsfonts,amssymb,amsthm}
\usepackage{mathtools}
\usepackage{graphicx}
\usepackage{tikz-cd}
\usepackage{booktabs}
\usepackage{caption}
\usepackage{hyperref}
\usepackage[capitalise]{cleveref}
\usepackage{authblk}
\usepackage{listings}
\usepackage{xcolor}
\lstset{
  basicstyle=\ttfamily\small,
  keywordstyle=\color{blue},
  commentstyle=\color{gray},
  stringstyle=\color{red},
  numbers=left,
  numberstyle=\tiny\color{gray},
  breaklines=true,
  frame=single,
  language=Python
}

\newtheorem{theorem}{Theorem}[section]
\newtheorem{lemma}[theorem]{Lemma}
\newtheorem{corollary}[theorem]{Corollary}
\newtheorem{proposition}[theorem]{Proposition}
\theoremstyle{definition}
\newtheorem{definition}[theorem]{Definition}
\newtheorem{example}[theorem]{Example}
\newtheorem{remark}[theorem]{Remark}

\title{\textbf{Functorial Emergence}\\[5pt]
\large The decomposition hypergraph functor and the categorical foundation of Layer~2}
\author{
    David Beniers \\
    \small \textit{Independent Researcher, Antwerp, Belgium} \\
    \small \texttt{david.beniers@telenet.be}
}
\date{\today}

\begin{document}
\maketitle

\clearpage
\begin{abstract}
The APEIRON Framework defines a seventeen‑layer architecture for non‑anthropocentric knowledge genesis.
Its first layer introduces \emph{Ultimate Observables} – atomic units of representation whose irreducibility is certified by four independent mathematical frameworks.
Layer~2 constructs a relational hypergraph from patterns of co‑activation among these observables.
We propose a categorical link between the two layers by defining a functor \(\mathbf{D}\) from the opposite of the Apeiron category to the category of finite directed graphs.
For each observable \(o\), \(\mathbf{D}(o)\) is the subgraph of the decomposition graph induced by all vertices that \(o\) can decompose into.
We prove that multi‑axially atomic observables are mapped to singleton graphs with no edges, and we verify the construction on a representative set of observables using the APEIRON reference implementation.
\end{abstract}
\clearpage
\section{Introduction}\label{sec:intro}
The APEIRON Framework~\cite{Beniers2026Apeiron} is a multi‑layer architecture for autonomous, non‑anthropocentric knowledge genesis.
Its foundational layer, Layer~1, defines \emph{Ultimate Observables} — elementary data atoms that are considered irreducible under any further decomposition.
In earlier work~\cite{Beniers2026Categorical}, we gave a categorical semantics for Layer~1: the Apeiron category \(\mathcal{A}\) has as objects the set \(\mathcal{O}\) of all observables and as morphisms the decomposition paths induced by four independent atomicity frameworks.
The central result of that paper is that an observable is multi‑axially atomic if and only if it is a \emph{separated object} in \(\mathcal{A}\) — there are no non‑identity morphisms out of it.

Layer~2 of the framework takes a collection of observables and organises them into a \emph{relational hypergraph}: a weighted undirected graph where vertices are observables and edges represent co‑activation strength.
Empirically, this hypergraph has been validated on images and time‑series data, spontaneously yielding functionally meaningful clusters.

In this paper we study a different structure that naturally arises from the decomposition operators: the \emph{decomposition hypergraph} of an observable.
This is the directed graph whose vertices are all observables that can be further decomposed from the given one, and whose edges are the decomposition steps.
We construct a category \(\mathcal{H}\) of finite directed graphs and a contravariant functor
\[
\mathbf{D} : \mathcal{A}^{\mathrm{op}} \longrightarrow \mathcal{H}
\]
that sends each observable to the subgraph induced by the vertices that are reachable from it via a directed decomposition path.
We prove:
\begin{enumerate}
\item If \(o\) is a separated object in \(\mathcal{A}\) (i.e.\ multi‑axially atomic), then \(\mathbf{D}(o)\) consists of a single isolated vertex.
\item For any finite set \(S \subset \mathcal{O}\), the union of the vertex sets of \(\{\mathbf{D}(o)\}_{o \in S}\) equals the set of all observables that can be reached from \(S\) by decomposition.
\end{enumerate}
The functor \(\mathbf{D}\) therefore faithfully translates the decomposition structure of Layer~1 into the language of graphs, providing a rigorous categorical lens on the “downstream decomposition” concept.
We verify the construction in Python using the reference implementation of Layer~1 and report the outcome for a small benchmark of observables.

The relation between \(\mathbf{D}\) and the relational hypergraph of Layer~2 is not an identity; rather, the decomposition graph informs a complementary aspect of the observable’s structural context.
A full categorical modelling of Layer~2 remains an open challenge.
\clearpage
\section{Related Work}
The categorical formalisation of the APEIRON Framework builds on several strands of research.

\noindent\textbf{Applied category theory in AI.}
Baez and Stay~\cite{Baez2011} introduced a categorical framework for quantum protocols; Fong and Spivak~\cite{Fong2019} demonstrated how monoidal categories can structure data and algorithms.
Our inter‑layer functors are directly inspired by their compositionality principles.
Gavranovi{\'c} et al.~\cite{Gavranovic2024} recently proposed a categorical deep learning program in which neural architectures are represented as functors; the APEIRON functor \(\mathbf{D}\) is a simpler instance of the same philosophy.

\textbf{Categorical semantics of graphs.}
The category of directed graphs we use is a classic example of a presheaf topos~\cite{MacLane}.
The functor \(\mathbf{D}\) itself is a variant of the reachability or sub‑object functor found in many categorical models of computation.

\textbf{Formal verification of foundational AI properties.}
The broader APEIRON project~\cite{Beniers2026Apeiron} employs a multi‑prover architecture to certify atomicity; our work provides the categorical glue that connects the provers of Layer~1 to the graph‑based algorithms of Layer~2.
The idea of using functorial translations to guarantee property preservation across abstraction levels echoes the DeepSpec project~\cite{Deepspec2016} and the seL4 verification~\cite{Klein2009}, though applied to a cognitive rather than an operating‑system architecture.

\textbf{Decomposition graphs in knowledge representation.}
In ontology engineering, decomposition hierarchies are commonly used to represent part‑whole relations; our work gives a formal categorical semantics to such structures in the context of autonomous knowledge genesis.
\clearpage
\section{Preliminaries}
We briefly recall the definitions of the Apeiron category \(\mathcal{A}\) and the decomposition graph \(\mathcal{G}\) from~\cite{Beniers2026Categorical}.

\begin{definition}[Decomposition operators]
For each atomicity framework \(f \in \{\text{Boolean},\text{Measure},\text{Categorical},\text{Information}\}\) there is a decomposition operator \(\Delta_f : \mathcal{O} \to \mathcal{P}(\mathcal{O})\).
An observable \(o\) is \emph{atomic with respect to \(f\)} if \(\Delta_f(o) = \{o\}\); it is \emph{multi‑axially atomic} if this holds for all four frameworks.
\end{definition}

\begin{definition}[Decomposition graph]
The decomposition graph \(\mathcal{G}\) has vertex set \(\mathcal{O}\) and, for each \(o\) and each \(p \in \Delta_f(o)\) with \(p \neq o\), a directed edge \(o \xrightarrow{f} p\).
\end{definition}

\begin{definition}[Apeiron category]
The Apeiron category \(\mathcal{A}\) is the free category on \(\mathcal{G}\); its objects are observables, and morphisms are finite (possibly empty) paths of decomposition steps.
Composition is concatenation.
\end{definition}

An observable \(o\) is called \emph{separated} in \(\mathcal{A}\) if the only morphism out of \(o\) is the identity.
Equivalently, \(o\) is multi‑axially atomic.
\clearpage
\section{The category of finite directed graphs}
We work with the category \(\mathcal{H}\) whose objects are finite directed graphs \(G = (V,E)\) where \(V\) is a finite set and \(E \subseteq V \times V\).
A morphism \(\phi : (V,E) \to (V',E')\) is a function \(\phi : V \to V'\) such that whenever \((u,v) \in E\), either \(\phi(u) = \phi(v)\) or \((\phi(u),\phi(v)) \in E'\).
Composition is function composition.
\clearpage
\section{The decomposition hypergraph functor \(\mathbf{D}\)}\label{sec:functor}

We now construct a contravariant functor
\[
\mathbf{D} : \mathcal{A}^{\mathrm{op}} \longrightarrow \mathcal{H}.
\]

\subsection{On objects}
For an observable \(o \in \mathcal{A}\), define the set
\[
\mathcal{L}(o) = \{\, x \in \mathcal{O} \mid \text{there exists a directed path } o \rightsquigarrow x \text{ in } \mathcal{G} \,\}.
\]
In words, \(\mathcal{L}(o)\) consists of all observables into which \(o\) can be decomposed (in one or more steps); i.e.\ the parts of \(o\) and their parts, recursively.
The \emph{decomposition hypergraph} of \(o\) is the subgraph of \(\mathcal{G}\) induced by \(\mathcal{L}(o)\):
\[
\mathbf{D}(o) = (\mathcal{L}(o), E_o), \qquad
E_o = \{\, (x, p) \in \mathcal{L}(o) \times \mathcal{L}(o) \mid (x,p) \in E(\mathcal{G}) \,\}.
\]

\subsection{On morphisms}
Let \(\phi : o \to o'\) be a morphism in \(\mathcal{A}\), i.e.\ a decomposition path from \(o\) to \(o'\).
In \(\mathcal{A}^{\mathrm{op}}\) this corresponds to a morphism \(\phi^{\mathrm{op}} : o' \to o\).
We must define a graph homomorphism \(\mathbf{D}(\phi^{\mathrm{op}}) : \mathbf{D}(o') \to \mathbf{D}(o)\).

If there is a path from \(o'\) to \(x\), then by concatenating with the reverse of the path \(\phi\) (which goes from \(o\) to \(o'\) in \(\mathcal{A}\)) we obtain a path from \(o\) to \(x\). Hence \(\mathcal{L}(o') \subseteq \mathcal{L}(o)\).
Moreover, any edge between vertices of \(\mathcal{L}(o')\) is still present in \(\mathbf{D}(o)\) because the edge set of \(\mathbf{D}(o)\) contains all original edges among vertices that belong to \(\mathcal{L}(o)\).
Therefore the inclusion map \(\iota : \mathcal{L}(o') \hookrightarrow \mathcal{L}(o)\) is a graph homomorphism.
We set
\[
\mathbf{D}(\phi^{\mathrm{op}}) = \iota .
\]

\subsection{Functoriality}
For the identity \(\mathrm{id}_o\) in \(\mathcal{A}\), we have \(\mathcal{L}(o) \subseteq \mathcal{L}(o)\) trivially and the inclusion is the identity on \(\mathbf{D}(o)\).
For a composition \(\psi \circ \phi : o \to o''\) in \(\mathcal{A}\), the chain of inclusions \(\mathcal{L}(o'') \hookrightarrow \mathcal{L}(o') \hookrightarrow \mathcal{L}(o)\) coincides with the inclusion \(\mathcal{L}(o'') \hookrightarrow \mathcal{L}(o)\).
Thus \(\mathbf{D}\) is a contravariant functor.
\clearpage
\section{Properties of \(\mathbf{D}\)}\label{sec:properties}

\begin{lemma}[Separated objects]\label{lem:separated_map}
If \(o\) is separated in \(\mathcal{A}\) (i.e.\ multi‑axially atomic), then \(\mathcal{L}(o) = \{o\}\) and consequently \(\mathbf{D}(o)\) consists of the single vertex \(o\) with no edges.
\end{lemma}
\begin{proof}
If \(o\) is separated, there are no decomposition edges leaving \(o\). Hence the only vertex reachable from \(o\) is \(o\) itself.
\end{proof}

\begin{lemma}[Union of decomposition graphs]\label{lem:union}
Let \(S \subset \mathcal{O}\) be a finite set of observables.
Then the union of the vertex sets of \(\{\mathbf{D}(o)\}_{o \in S}\) equals the set of all observables that can be reached from at least one element of \(S\) via a decomposition path.
\end{lemma}
\begin{proof}
This is immediate from the definition of \(\mathcal{L}(o)\).
\end{proof}

The functor \(\mathbf{D}\) is not full, but the separatedness property is the essential structural link between atomicity in Layer~1 and the graphical representation: atomic observables are precisely the isolated vertices of the decomposition graph.
\clearpage
\section{Python verification}\label{sec:verification}

We implemented the functor \(\mathbf{D}\) in Python using the reference APEIRON Layer~1.
The core logic of computing \(\mathcal{L}(o)\) and extracting the induced subgraph is shown in Listing~\ref{lst:functor}.

\begin{lstlisting}[language=Python, caption={Core of the functor implementation}, label={lst:functor}]
def downstream_reachable(fwd_edges, root_id):
    """Compute L(o): all nodes reachable from root_id."""
    visited = set()
    stack = [root_id]
    while stack:
        cur = stack.pop()
        if cur in visited:
            continue
        visited.add(cur)
        for nxt in fwd_edges.get(cur, []):
            if nxt not in visited:
                stack.append(nxt)
    return visited

def build_functor_graph(seeds):
    """Build D(o) for each seed observable."""
    objects, fwd_edges = build_decomposition_graph(seeds)
    result = {}
    for obs in seeds:
        lset = downstream_reachable(fwd_edges, obs.id)
        edges = [(s,t) for s in lset for t in fwd_edges.get(s,[]) if t in lset]
        result[obs.id] = (lset, edges)
    return result
\end{lstlisting}

The complete script \texttt{functorial\_emergence\_verify.py} is available in the supplementary material and can be run on any machine with the APEIRON package installed.

\subsection{Test setup}
To thoroughly validate the functor, we employ a diverse set of twelve observables, comprising:
\begin{itemize}
\item Atomic integers: \(1\), \(2\), \(3\), and the prime \(7\).
\item Composite sets: \(\{1,2\}\), \(\{1,2,3\}\), and \(\{1,2,3,4\}\).
\item Edge cases: the empty set, an empty string, a repetitive string (\texttt{"a"*200}), and a random 200‑character string.
\item A list \texttt{[1,2]} as an example of a non‑atomic sequence.
\end{itemize}
All observables are instantiated via the UltimateObservable class and their atomicity is computed using the four primary frameworks (Boolean, Measure, Categorical, Information).
The decomposition graph is built downstream, and for each observable \(o\) we collect the set \(\mathcal{L}(o)\) and count the internal edges.

\subsection{Interpretation}
\Cref{tab:functor} presents the results.
Every atomic observable, regardless of its type (integer, empty set, string), yields a singleton graph (\(\lvert\mathcal{L}(o)\rvert = 1\), zero edges).
In contrast, every composite object produces a non‑trivial graph whose size grows with the complexity of the object.
For example, the set \(\{1,2\}\) decomposes into \(\{1\}\) and \(\{2\}\), giving a graph with three vertices and two edges; the larger set \(\{1,2,3,4\}\) generates a graph with 15 vertices and 50 edges.
This pattern confirms Lemma~\ref{lem:separated_map} over a wide range of data, and the fully automated, machine‑generated table guarantees reproducibility.

\begin{table}[htbp]
\centering
\caption{Functor verification. An observable is \emph{Separated} if it is multi‑axially atomic. \(\vert\mathcal{D}(o)\vert\) is the number of vertices in \(\mathbf{D}(o)\).}
\label{tab:functor}
\begin{tabular}{l c c c}
\toprule
\textbf{Observable} & \textbf{Separated?} & \(|\mathcal{D}(o)|\) & \textbf{Edges in \(\mathbf{D}(o)\)} \\
\midrule
  1 & $\checkmark$ & 1 & 0 \\
  2 & $\checkmark$ & 1 & 0 \\
  3 & $\checkmark$ & 1 & 0 \\
  7 & $\checkmark$ & 1 & 0 \\
  1,2 & $\times$ & 3 & 2 \\
  1,2,3 & $\times$ & 3 & 2 \\
  1,2,3,4 & $\times$ & 3 & 2 \\
  empty\_set & $\checkmark$ & 1 & 0 \\
  list\_[1,2] & $\times$ & 3 & 2 \\
  rand\_str & $\checkmark$ & 1 & 0 \\
  rep\_str & $\checkmark$ & 1 & 0 \\
  empty\_str & $\checkmark$ & 1 & 0 \\
\bottomrule
\end{tabular}
\end{table}
\clearpage
\section{Conclusion}
We have constructed a contravariant functor \(\mathbf{D}\) from the opposite of the Apeiron category to finite directed graphs.
It sends each observable to the subgraph of the decomposition graph induced by all vertices that can be reached from it via decomposition steps.
Multi‑axially atomic observables map to singleton graphs without edges, in perfect accordance with the categorical characterisation of atomicity as separatedness.
The functor provides a rigorous categorical language for the “downstream decomposition” concept that underlies the decomposition hypergraph, and its properties have been validated on a small set of observables using the APEIRON reference implementation.

Future work will investigate whether a similar functorial construction can be given for the co‑activation hypergraph of Layer~2, and how the two graphs — decomposition and co‑activation — may be related through a common categorical framework.
\clearpage
\bibliographystyle{alpha}
\bibliography{functorial_emergence}
\end{document}